\section{Proofs for the box-face characterization}
\label{app:face}

\subsection{A per-face lower bound}
\label{app:face:lower-bound}

Let $F$ be a root face with free-coordinate set $J$ and dimension
$d\ge1$.  If a test box $B=\prod_i[a_i,b_i]\subseteq D$ meets $F$,
each coordinate fixed on $F$ is also an endpoint of the corresponding
interval of $B$.  Its gap factor therefore vanishes on $F\cap B$.
For the free coordinates, \eqref{face:eq:validity} and the
arithmetic--geometric mean inequality give, almost everywhere on the
relative interior of a nondegenerate $F\cap B$,
\[
 (m+\varepsilon)^{-d/2}
 \le d^{-d/2}\left(\prod_{i\in J}\alpha_i\right)^{-1/2}
       \prod_{i\in J}\bigl((x_i-a_i)(b_i-x_i)\bigr)^{-1/2}.
\]
Each one-dimensional factor has integral
\[
 \int_{a_i}^{b_i}\frac{dt}{\sqrt{(t-a_i)(b_i-t)}}=\pi.
\]
It follows that
\begin{equation}
 \label{face:eq:one-box-integral}
 \int_{F\cap B}(m+\varepsilon)^{-d/2}\,d\sigma_F
 \le\pi^d d^{-d/2}
       \left(\prod_{i\in J}\alpha_i\right)^{-1/2}.
\end{equation}
If a free-coordinate interval has zero width, the intersection has
zero $d$-measure and contributes zero instead.  Summing
\eqref{face:eq:one-box-integral} over any valid cover proves
\eqref{face:eq:lower-explicit}.  The same argument integrates over
$F\cap A_e$ for an owned set $A_e\subseteq B_e$ and then bounds it
by the product integral over $F\cap B_e$; full validity on the
unowned boundary is unnecessary.

\subsection{Projection and relative sublevel cubes}
\label{app:face:projection}

We first work on a cube $D=\prod_i[L_i,L_i+s_0]$.
Choose a positive Lipschitz constant $M$ for $\nabla m$, increasing
it if necessary.  Convexity of $D$ gives the two-sided Taylor estimate
\begin{equation}
 \label{face:eq:descent}
 \left|m(z)-m(y)-\nabla m(y)^\top(z-y)\right|
 \le\frac M2\|z-y\|_2^2\qquad(y,z\in D).
\end{equation}
Fix $0<s\le s_0/2$ and $y\in D$.  Move each coordinate whose
distance from a root endpoint is strictly less than $s$ to that
endpoint.  If both endpoints qualify at $s=s_0/2$, choose either;
under the strict inequality they cannot both qualify.  Let $F$ be
the resulting face, let $d=\dim F$, and let $y'$ be the moved point.
Every free coordinate of $y'$ has distance at least $s$ from both
endpoints.  In particular, the relative cube
\[
 Q_F(y',s)=\{z\in\operatorname{aff}F:
                 \|z-y'\|_\infty\le s\}
\]
is contained in $F$.

For a moved coordinate $i$, let $\eta_i$ be the sign pointing from
$y_i$ to its chosen endpoint.  There is room for the step
$y-\eta_i s e_i$ in the opposite direction.  Nonnegativity and
\eqref{face:eq:descent} imply
\[
 0\le m(y-\eta_i s e_i)
 \le m(y)-s\eta_i\partial_i m(y)+\frac M2s^2,
 \qquad
 \eta_i\partial_i m(y)\le\frac{m(y)}s+\frac M2s.
\]
Coordinates already at their chosen endpoint have zero displacement
and require no derivative test.  Summing the first-order terms for
the other moved coordinates and using \eqref{face:eq:descent} gives
\begin{equation}
 \label{face:eq:projection-height}
 m(y')\le(1+n-d)m(y)+(n-d)Ms^2.
\end{equation}
For each free coordinate, both steps $y'\pm se_i$ are feasible.
The same argument gives
$|\partial_i m(y')|\le m(y')/s+Ms/2$.
For every $z\in Q_F(y',s)$,
\begin{align}
 m(z)&\le(1+d)m(y')+dMs^2\notag\\
 &\le(1+d)(1+n-d)m(y)
       +\bigl((1+d)(n-d)+d\bigr)Ms^2\notag\\
 &\le K_n\bigl(m(y)+Ms^2\bigr),
 \qquad K_n=\frac{(n+2)^2}{4}.
 \label{face:eq:projection-fatness}
\end{align}
Both coefficients in the preceding line are at most $K_n$.
These inequalities use nonnegativity only on the closed cube.

\subsection{Dyadic packing and the vertex estimate}
\label{app:face:characterization}

Let $s_j=s_0 2^{-j}$ and $\bar\alpha=\max_i\alpha_i$.
Put
\[
 \Lambda_0=\bar\alpha n/4,\qquad
 \Lambda_1=K_n(\Lambda_0+M),\qquad
 \Lambda_2=\Lambda_0+\Lambda_1.
\]
An invalid level-$j$ cube $B$ has a point $y\in B$ with
\begin{equation}
 \label{face:eq:dyadic-witness}
 m(y)+\varepsilon
 <\sum_i\alpha_i(y_i-a_i)(b_i-y_i)
 \le\Lambda_0s_j^2.
\end{equation}
Thus no box remains invalid once
$\Lambda_0s_j^2\le\varepsilon$, and the dyadic tree is finite.
For $j\ge1$, apply \eqref{face:eq:projection-fatness} to the witness
with $s=s_j$.  It gives a face $F$ and a point $y'$ at sup-distance
less than $s_j$ from $y$, with
\[
 Q_F(y',s_j)\subseteq F\cap\{m\le\Lambda_1s_j^2\}.
\]

For each face, choose a maximal subset of these projected witnesses
whose mutual sup-distances exceed $s_j$.  This is a finite selection
from the invalid boxes at that level.  On a face of dimension $d\ge1$,
the relative cubes of radius $s_j/2$ about the selected points have
disjoint interiors, measure $s_j^d$, and lie in the displayed
sublevel set.  Hence the number selected is at most
\[
 s_j^{-d}\sigma_F\{m\le\Lambda_1s_j^2\}.
\]
By maximality, every projected witness is within $s_j$ of a selected
point.  Its original witness is within $2s_j$, and its entire grid
box is within $3s_j$.  A cube of radius $3s_j$ meets at most $8^n$
level-$j$ grid boxes.  Consequently the number of invalid boxes
charged to positive-dimensional faces is at most
\begin{equation}
 \label{face:eq:level-packing}
 8^n\sum_{F\in\mathcal F_+(D)}
 s_j^{-d_F}\sigma_F\{m\le\Lambda_1s_j^2\}.
\end{equation}
Shared grid boundaries only decrease the useful packing volume and
are covered by this deliberately loose grid-count constant.

Whenever a level contributes, \eqref{face:eq:dyadic-witness} implies
$\varepsilon<\Lambda_0s_j^2$, so the sublevel in
\eqref{face:eq:level-packing} lies in
$\{m+\varepsilon\le\Lambda_2s_j^2\}$.
For $t>0$ and $d\ge1$, the geometric series gives
\[
 \sum_{j\ge1}s_j^{-d}\mathbf1\{t\le\Lambda_2s_j^2\}
 \le\frac{1}{1-2^{-d}}\left(\frac{\Lambda_2}t\right)^{d/2}
 \le2\Lambda_2^{d/2}t^{-d/2}.
\]
Tonelli's theorem therefore bounds the sum over levels in
\eqref{face:eq:level-packing} by
\begin{equation}
 \label{face:eq:packing-sum}
 2\cdot8^n\sum_{F\in\mathcal F_+(D)}
 \Lambda_2^{d_F/2}I_F(\varepsilon).
\end{equation}

If the projected face is a vertex $v$, the witness is strictly within
$s_j$ of $v$ in every coordinate.  It is therefore in the unique
corner grid box at that vertex.  Let $V_v(\varepsilon)$ count the
levels $j\ge1$ at which that corner box is invalid.  The projection
estimate also applies when all coordinates of a corner-box witness
are moved to $v$ with distances at most $s_j$: the opposite
length-$s_j$ steps stay in $D$ because $2s_j\le s_0$.
It gives
$m(v)\le(1+n)m(y)+nMs_j^2\le\Lambda_1s_j^2$ on every such level.
For $m(v)>0$, this
allows only finitely many levels, independently of
$\varepsilon$.

It remains to control $V_v$ when $m(v)=0$.  Use inward edge
coordinates $t_i\ge0$ at $v$ and write $g_i\ge0$ for the inward
derivatives.  Put $g=\min_i g_i$ and choose an edge $E$ attaining
this minimum.  On a corner box of side $s$,
\[
 m(v+t)\ge g\|t\|_1-\frac M2\|t\|_2^2
          \ge\left(g-\frac M2s\right)\|t\|_1,
 \qquad
 \sum_i\alpha_i t_i(s-t_i)\le\bar\alpha s\|t\|_1.
\]
The box is valid whenever
$s\le g/(\bar\alpha+M/2)$.  Independently, it is valid whenever
$\Lambda_0s^2\le\varepsilon$.  Define
\[
 \rho=\max\{g/M,\sqrt\varepsilon\},\qquad
 c_* = \min\{M/(\bar\alpha+M/2),\Lambda_0^{-1/2}\}>0.
\]
An invalid corner box must have
$s>\max\{g/(\bar\alpha+M/2),\sqrt{\varepsilon/\Lambda_0}\}
\ge c_*\rho$.  Counting dyadic scales, with
$\log_+t=\max\{0,\log t\}$, gives
\begin{equation}
 \label{face:eq:vertex-levels}
 V_v(\varepsilon)
 \le1+\frac{\log_+(1/c_*)}{\log2}
       +\frac{\log_+(s_0/\rho)}{\log2}.
\end{equation}
Along $E$, Taylor's estimate gives
$m(v+te)\le gt+Mt^2/2$.
For $\rho\le t\le s_0$ this is at most $3Mt^2/2$ and
$\varepsilon\le t^2$.  It follows, including the case
$\rho>s_0$ where the right side is zero, that
\begin{equation}
 \label{face:eq:edge-absorbs-vertex}
 I_E(\varepsilon)
 \ge(3M/2+1)^{-1/2}\log_+(s_0/\rho).
\end{equation}
Equations \eqref{face:eq:vertex-levels} and
\eqref{face:eq:edge-absorbs-vertex} absorb the vertex count into an
edge integral and a fixed constant.  There are finitely many vertices,
and each edge can be selected by at most its two endpoints.

Every invalid dyadic node has exactly $2^n$ children.  If $H$ is the
number of invalid nodes, the total number tested is $1+2^nH$.
The root contributes at most one invalid node.  Combining this identity,
\eqref{face:eq:packing-sum}, the fixed contribution of nonoptimal
vertices, and the optimal-vertex estimate proves
\eqref{face:eq:upper-explicit} on a cube.  Replacing each simultaneous
coordinate refinement by $n$ successive binary coordinate splits gives
the same terminal partition.  The intermediate binary nodes add at most
a fixed multiple of the dyadic work.  This proves the tree claim.

For a general box, set $w_i=u_i-\ell_i$ and use
$x_i=\ell_i+w_i z_i$ to map $[0,1]^n$ onto $D$.
The transformed coefficients are $\alpha_iw_i^2>0$, and the transformed
gradient is Lipschitz.  On each face,
\[
 I_F(\varepsilon)=
 \left(\prod_{i\in J_F}w_i\right)
 \int_{\widetilde F}(\widetilde m+\varepsilon)^{-d_F/2}
      \,d\sigma_{\widetilde F}.
\]
These are fixed positive factors.  The transformation preserves covers,
rectangular partitions, and coordinate split trees, and maps the unit
dyadic tree to coordinate-wise halving on $D$.
The upper bound therefore holds on $D$ with changed constants.
The per-face lower bound in \cref{app:face:lower-bound} was already
proved in the original coordinates.  A finite sum and its maximum are
comparable.  Together with
\eqref{face:eq:benchmarks}, these bounds prove all the comparisons in
\cref{face:thm:characterization}.

\subsection{Dominance by the full-box integral}
\label{app:face:full-box}

We prove case (ii) of \cref{face:prop:full-box}.  By the preceding
diagonal change of coordinates it suffices to use $D=[0,1]^n$.
Choose $r>0$ such that $m\ge0$ and $\nabla m$ has positive Lipschitz
constant $M$ on $D+\overline B_2(0,r)$.  For $z\in D$ with
$m(z)\le t\le Mr^2/2$, a descent step of length $r$ in the negative
gradient direction first gives
\[
 \|\nabla m(z)\|\le t/r+Mr/2\le Mr.
\]
If the gradient is zero this conclusion is immediate.
Thus the step $-\nabla m(z)/M$ is admissible, and the descent estimate
and nonnegativity give the sharper inequality
\begin{equation}
 \label{face:eq:gradient-nonnegative}
 \|\nabla m(z)\|^2\le2Mm(z)\le2Mt.
\end{equation}
Put $t_0=\min\{Mr^2/2,1/4\}$ and $K=1+\sqrt{2Mn}+Mn/2$.
For a proper face $F$ of dimension $d\ge1$, write
\[
 V_F(t)=\sigma_F\{z\in F:m(z)\le t\},\qquad
 V_D(t)=\vol\{z\in D:m(z)\le t\}.
\]
For $0<t\le t_0$, put $s=\sqrt t$ and select a maximal set of points
in $F\cap\{m\le t\}$ separated by more than $2s$ in sup norm.
Compactness gives a finite such set, say of size $P$.
Its relative cubes of radius $2s$ cover the sublevel set, so
$V_F(t)\le P(4s)^d$.

For each selected point $z$, choose an axis-parallel cube of side
$s$ contained in $D$ and having $z$ as a vertex.  Such a cube exists
because $s\le1/2$: in each free coordinate choose a direction with
at least $1/2$ of room, and in each fixed coordinate point inward.
Every point $x$ of this cube has distance at most $\sqrt n\,s$ from
$z$.  Equation \eqref{face:eq:gradient-nonnegative} and Taylor's
estimate give $m(x)\le Kt$.
The selected cubes have disjoint interiors, because their vertex
sup-distances exceed $2s$.  Consequently
\begin{equation}
 \label{face:eq:face-volume-transfer}
 V_D(Kt)\ge P s^n
 \ge4^{-d}t^{(n-d)/2}V_F(t)
 \qquad(0<t\le t_0).
\end{equation}

Let $a=d/2$, $b=n/2$, and assume $0<\varepsilon<t_0$.
The elementary layer-cake identity is
\[
 I_F(\varepsilon)
 =a\int_0^\infty(t+\varepsilon)^{-a-1}V_F(t)\,dt.
\]
The part $t\ge t_0$ is bounded by $\sigma_F(F)t_0^{-a}$.
On $[0,\varepsilon]$, monotonicity bounds the contribution by
$\varepsilon^{-a}V_F(\varepsilon)$.  By
\eqref{face:eq:face-volume-transfer} this is at most
$4^d\varepsilon^{-b}V_D(K\varepsilon)$, while
\[
 I_D(\varepsilon)
 \ge((K+1)\varepsilon)^{-b}V_D(K\varepsilon).
\]
On $[\varepsilon,t_0]$, the volume transfer and
$t\ge(t+\varepsilon)/2$ bound the integrand by a constant times
$(t+\varepsilon)^{-b-1}V_D(Kt)$.
Since $Kt+\varepsilon\le K(t+\varepsilon)$, this is at most a
constant times $(Kt+\varepsilon)^{-b-1}V_D(Kt)$.
Substituting $u=Kt$ and using the layer-cake identity for $I_D$
bounds its integral by $C I_D(\varepsilon)$.
Thus $I_F(\varepsilon)\le C_0+C_1 I_D(\varepsilon)$ for every proper
positive-dimensional face.  This proves the claimed dominance and
completes \cref{face:prop:full-box}.
